\documentclass[11pt]{article}

\usepackage[margin=1in]{geometry}
\usepackage[T1]{fontenc}
\usepackage[utf8]{inputenc}
\usepackage{lmodern}
\usepackage{amsmath,amssymb}
\usepackage{booktabs}
\usepackage{array}
\usepackage{enumitem}
\usepackage{longtable}
\usepackage{xcolor}
\usepackage{hyperref}
\usepackage{listings}

\hypersetup{
    colorlinks=true,
    linkcolor=blue,
    urlcolor=blue
}

\lstset{
    basicstyle=\ttfamily\small,
    breaklines=true,
    columns=fullflexible,
    keepspaces=true,
    frame=single
}

\title{Static Branches}
\author{}
\date{}

\begin{document}
\maketitle

\section{Purpose}

The static side of the project now has two complementary branches:
\begin{itemize}[leftmargin=2em]
    \item \texttt{static\_artificial}, a clean synthetic baseline,
    \item \texttt{static\_campus\_area\_1}, a semantic image-based campus branch.
\end{itemize}

Both remain static in the strict sense that there are no time-varying obstacle schedules.

\section{\texttt{static\_artificial}}

This branch is still the cleanest place to compare the two mappings without image semantics. For each sampled run configuration, the code:
\begin{enumerate}[leftmargin=2em]
    \item creates a fresh artificial composite grid,
    \item places static obstacles according to the configured density,
    \item builds classical and cyclic mapped versions,
    \item samples dispersed one-to-one starts and dispersed one-to-one goals.
\end{enumerate}

Because the branch is synthetic, it is best read as the lowest-assumption baseline for the whole study.

\section{\texttt{static\_campus\_area\_1}}

This branch now uses the campus-area-1 image as a \emph{static} semantic map rather than a dynamic branch source. The campus loader keeps the meaning of the image colors:
\begin{itemize}[leftmargin=2em]
    \item zone colors are traversable and spawnable,
    \item white walkways are traversable but non-spawnable,
    \item gray regions are non-traversable.
\end{itemize}

The branch then samples:
\begin{itemize}[leftmargin=2em]
    \item a dispersed start set from one campus zone,
    \item a clustered goal set from the other campus zone, sampled according to the global \texttt{compact\_clustering} rule,
    \item a one-to-one assignment between those sets.
\end{itemize}

Because the two groups must occupy different zones, the branch keeps the intended pressure toward cross-zone travel and bottleneck usage even though the environment itself is static.

\section{Reachability and Fairness}

The semantic campus static branch uses individual reachability checks during assignment, while the artificial static branch keeps its older unrestricted dispersed assignment behavior. This difference is deliberate: the campus branch is meant to study zone-separated spatial structure, not to waste most samples on trivially disconnected start-goal pairs.

\section{Why Both Static Branches Matter}

Taken together, the two static branches answer different questions:
\begin{itemize}[leftmargin=2em]
    \item \texttt{static\_artificial}: what the mapping difference looks like in a controlled synthetic setting,
    \item \texttt{static\_campus\_area\_1}: what the mapping difference looks like on a semantically meaningful real image while still avoiding time variation.
\end{itemize}

\end{document}
